% Part: many-valued-logic
% Chapter: syntax-and-semantics
% Section: introduction

\documentclass[../../../include/open-logic-section]{subfiles}

\begin{document}

\olfileid{mvl}{syn}{int}

\olsection{પરિચય}

શાસ્ત્રીય તર્કશાસ્ત્રમાં આપણે $\lnot$, $\land$, $\lor$,
$\lif$ અને $\liff$ એ વિધાનાત્મક સંયોજકો વડે
!!{propositional variable}sમાંથી બનેલાં !!{formula}sનો
અભ્યાસ કરીએ છીએ. શાસ્ત્રીય તર્કશાસ્ત્રનું અર્થવિજ્ઞાન
વ્યાખ્યાયિત કરતી વખતે $\True$ અને $\False$ એ બે
સત્યમૂલ્યો વાપરીએ છીએ. !!a{valuation}~$\pAssign{v}$માં
!!{propositional variable}sને અર્થ આપીએ છીએ; તે
!!{propositional variable}sને $\True$, $\False$ એ
સત્યમૂલ્યો નિયુક્ત કરે છે. દરેક !!{valuation} પછી
કોઈ પણ !!{formula}~$!A$ માટે સત્યમૂલ્ય
$\pValue{v}(!A)$ નક્કી કરે છે. જ્યારે
$\pValue{v}(!A) = \True$ હોય ત્યારે અને માત્ર
ત્યારે !!^a{formula}, !!a{valuation}~$\pAssign{v}$માં
સંતોષાય છે; તેને $\pSat{v}{!A}$ લખીએ છીએ.

બહુમૂલ્ય તર્કશાસ્ત્રો $\True$ અને $\False$ ઉપરાંત
વધુ સત્યમૂલ્યો માન્ય રાખીને શાસ્ત્રીય દ્વિમૂલ્ય
તર્કશાસ્ત્રનું સામાન્યકરણ કરે છે. તેથી બહુમૂલ્ય
તર્કશાસ્ત્રમાં !!a{valuation}~$\pAssign{v}$ એ
દરેક !!{propositional variable}~$p$ને શક્ય
સત્યમૂલ્યોમાંથી એક નિયુક્ત કરતું વિધેય છે. માન્ય
સત્યમૂલ્યોના ગણને સામાન્ય રીતે~$V$ કહીશું.
શાસ્ત્રીય તર્કશાસ્ત્ર પણ $V = \{\True, \False\}$
વાળું બહુમૂલ્ય તર્કશાસ્ત્ર છે; તેમાં સંયોજકો માટેની
પરિચિત સત્યતાસારણીઓ વડે સત્યમૂલ્ય
$\pValue{v}(!A)$ ગણાય છે.

વધારાનાં સત્યમૂલ્યો ઉમેરીએ પછી ``અને'', ``અથવા'',
``નહિ'' અને ``જો---તો'' તરીકે વાંચાતા સંયોજકો માટે
$\pValue{v}(!A)$ કેવી રીતે ગણવું તેની એકથી વધુ
સ્વાભાવિક પસંદગીઓ મળે છે. તેથી બહુમૂલ્ય તર્કશાસ્ત્ર
માત્ર સત્યમૂલ્યોના ગણથી જ નહીં, પરંતુ દરેક
સંયોજક માટે પસંદ કરેલાં \emph{સત્યવિધેયો}થી પણ
નક્કી થાય છે. પરંતુ બહુમૂલ્ય તર્કશાસ્ત્ર~$\Log L$
માટે આ પસંદગી થઈ જાય પછી, શાસ્ત્રીય તર્કશાસ્ત્રની
જેમ જ, મૂલ્યાંકન સત્યમૂલ્ય
$\pValue{v}(!A)[\Log L]$ને એકમાત્ર રીતે નક્કી કરે
છે. શાસ્ત્રીય તર્કશાસ્ત્રની જેમ બહુમૂલ્ય તર્કશાસ્ત્રો
પણ \emph{સત્યતાફલનલક્ષી} છે.

અર્થવિજ્ઞાનના આ મૂળ ઘટકો મળ્યા પછી પુનરુક્તિ,
ફલિતતા અને સંતોષ્યતા જેવા અર્થવિજ્ઞાનના ખ્યાલોનો
અનુરૂપ અર્થ વ્યાખ્યાયિત કરી શકીએ. શાસ્ત્રીય
તર્કશાસ્ત્રમાં કોઈ પણ~$\pAssign{v}$ માટે સત્યમૂલ્ય
$\pValue{v}(!A) = \True$ હોય તો !!a{formula}
પુનરુક્તિ છે. બહુમૂલ્ય તર્કશાસ્ત્રમાં આને પણ થોડું
સામાન્ય બનાવવું પડે છે. સૌપ્રથમ, સત્યમૂલ્યોના
ગણ~$V$માં $\True$ હોવું જરૂરી નથી. દાખલા તરીકે,
કેટલાક બહુમૂલ્ય તર્કશાસ્ત્રો સત્યમૂલ્યોના ગણ તરીકે
$0$ અને~$1$ વચ્ચેની બધી પરિમેય સંખ્યાઓ જેવી
સંખ્યાઓ વાપરે છે. એવા કિસ્સામાં $1$ સામાન્ય રીતે
$\True$ની ભૂમિકા ભજવે છે. અન્ય તર્કશાસ્ત્રોમાં
માત્ર એક નહીં, અનેક સત્યમૂલ્યો એવી ભૂમિકા
ભજવે છે. તેથી દરેક બહુમૂલ્ય તર્કશાસ્ત્રમાં
\emph{નિર્દિષ્ટ મૂલ્યો}નો ગણ~$V^+$ હોય એવી શરત
રાખીએ છીએ. પછી $\pValue{v}(!A)[\Log L] \in V^+$
હોય ત્યારે અને માત્ર ત્યારે !!a{formula},
!!a{valuation}~$\pAssign{v}$માં સંતોષાય છે;
તેને $\pSat{v}{!A}[\Log L]$ લખીએ છીએ.
દરેક~$\pAssign{v}$ માટે $\pValue{v}(!A) \in
V^+$ હોય તો અને તો જ !!^a{formula}~$!A$ એ
તર્કશાસ્ત્રની પુનરુક્તિ છે, જેને $\Entails[\Log L]
!A$ લખીએ છીએ. અંતે, !!{formula}sના ગણ~$\Gamma$માંથી
$!A$ ફલિત થાય છે એમ કહીએ અને $\Gamma
\Entails[\Log L] !A$ લખીએ છીએ, જો
તે ગણનાં બધાં !!{formula}sને સંતોષતું દરેક
!!{valuation}~$!A$ને પણ સંતોષે.

\end{document}
